% Generated by tools/audit_to_tex.py from the audit note. Do not edit by hand.
\section{Basis of the statements}\label{app:audit}

This appendix records the basis of the paper's statements about the proof. \textsf{M}: computed from the compiled archive. \textsf{S}: the Lean statement was read and transcribed by the authors. \textsf{P}: the Lean proof was read by the authors. \textsf{C}: the Lean statement was compared by the authors with the published theorem as they know it, without opening the papers. \textsf{R}: taken from a reader's report, published with this paper. Paths are relative to the archive's \lean{src} directory, and a line number follows a colon; paths beginning \lean{results/} are in \lean{verify/} of the repository of \cref{app:reproduce}, and those beginning \lean{refs/} beside the source of this paper.

{\footnotesize\setlength{\tabcolsep}{4pt}
\begin{longtable}{@{}>{\raggedright\arraybackslash}p{0.36\textwidth}>{\raggedright\arraybackslash}p{0.12\textwidth}>{\raggedright\arraybackslash}p{0.46\textwidth}@{}}
\toprule
Claim & Level & What was read\\
\midrule
\endhead
\bottomrule
\endfoot
The archive's size and hash, the counts of files and lines, and the static inspection & \textsf{M} & \leanb{results/\allowbreak{}scan.\allowbreak{}txt}\\
\addlinespace[2pt]
The five theorems of \leanb{Main.\allowbreak{}lean} and their statements & \textsf{S} & \leanb{Main.\allowbreak{}lean:\allowbreak{}4-\allowbreak{}12}; \leanb{Statements.\allowbreak{}lean}\\
\addlinespace[2pt]
\leanb{Riemann\allowbreak{}Hypothesis} and \leanb{riemann\allowbreak{}Zeta} are Mathlib's & \textsf{M} & \leanb{results/\allowbreak{}check.\allowbreak{}txt}\\
\addlinespace[2pt]
The build succeeds; every module is built & \textsf{M} & \leanb{results/\allowbreak{}build-\allowbreak{}summary.\allowbreak{}txt}\\
\addlinespace[2pt]
Each of the five theorems depends on \leanb{propext}, \leanb{Classical.\allowbreak{}choice}, \leanb{Quot.\allowbreak{}sound} only & \textsf{M} & \leanb{results/\allowbreak{}axioms.\allowbreak{}txt}\\
\addlinespace[2pt]
The dependency cone and every statement that a declaration is or is not used by the proof & \textsf{M} & \leanb{results/\allowbreak{}cut-\allowbreak{}result.\allowbreak{}txt}, \leanb{results/\allowbreak{}cut-\allowbreak{}in.\allowbreak{}txt.\allowbreak{}xz}, \leanb{results/\allowbreak{}cut-\allowbreak{}out.\allowbreak{}txt.\allowbreak{}xz}, \leanb{results/\allowbreak{}cut-\allowbreak{}modules.\allowbreak{}txt.\allowbreak{}xz}, \leanb{results/\allowbreak{}cut-\allowbreak{}five-\allowbreak{}extra.\allowbreak{}txt}\\
\addlinespace[2pt]
The kernel replay and its counts & \textsf{M} & \leanb{results/\allowbreak{}replay-\allowbreak{}summary.\allowbreak{}txt}, \leanb{results/\allowbreak{}replay-\allowbreak{}axioms.\allowbreak{}txt}, \leanb{results/\allowbreak{}replay-\allowbreak{}skipped.\allowbreak{}txt}, \leanb{results/\allowbreak{}replay-\allowbreak{}log.\allowbreak{}txt}\\
\addlinespace[2pt]
What the elaborators of the archive do & \textsf{P} & \leanb{Formalism/\allowbreak{}Hermitian\allowbreak{}Curvature/\allowbreak{}Hermitian\allowbreak{}Curvature\allowbreak{}Proofs/\allowbreak{}Cut\allowbreak{}Cpt\allowbreak{}Copy.\allowbreak{}lean:\allowbreak{}87}; \leanb{Formalism/\allowbreak{}Local\allowbreak{}To\allowbreak{}Global/\allowbreak{}Indefinite\allowbreak{}Real/\allowbreak{}Hecke\allowbreak{}Nonvanishing/\allowbreak{}Imports.\allowbreak{}lean:\allowbreak{}5}; \leanb{Formalism/\allowbreak{}Local\allowbreak{}To\allowbreak{}Global/\allowbreak{}Number\allowbreak{}Field\allowbreak{}Families/\allowbreak{}Own\allowbreak{}Struct.\allowbreak{}lean:\allowbreak{}53-\allowbreak{}60}; \leanb{Build\allowbreak{}Aid/\allowbreak{}D1.\allowbreak{}lean} to \leanb{D5.\allowbreak{}lean}\\
\addlinespace[2pt]
\leanb{riemann\allowbreak{}Hypothesis} is \leanb{rh\allowbreak{}Of\allowbreak{}All}, which is \leanb{Formalism.\allowbreak{}VW.\allowbreak{}rh\_\allowbreak{}of\_\allowbreak{}every all\allowbreak{}Chi} & \textsf{P} & \leanb{Main.\allowbreak{}lean:\allowbreak{}4}; \leanb{Supplements/\allowbreak{}Kernels/\allowbreak{}Main\allowbreak{}Theorems.\allowbreak{}lean:\allowbreak{}9}\\
\addlinespace[2pt]
The three hypotheses \leanb{Green}, \leanb{VWin}, \leanb{Nonempty (Hyps\allowbreak{}Alt \ldots{})}; the window exponent 3/2; hosts from \leanb{supply\allowbreak{}Chi} & \textsf{P} & \leanb{Development/\allowbreak{}Automorphic\allowbreak{}Arithmetic/\allowbreak{}Theta\allowbreak{}Correspondence/\allowbreak{}Theta\allowbreak{}Pair\allowbreak{}Lifts/\allowbreak{}Third\allowbreak{}Place\allowbreak{}Series/\allowbreak{}Top\allowbreak{}Plug.\allowbreak{}lean:\allowbreak{}12-\allowbreak{}18}\\
\addlinespace[2pt]
The assembly of the record's fields into the links and the final call & \textsf{P} & \leanb{\ldots{}/\allowbreak{}Third\allowbreak{}Place\allowbreak{}Series/\allowbreak{}Top\allowbreak{}Shape.\allowbreak{}lean}, whole file\\
\addlinespace[2pt]
The definition of G and of a family & \textsf{S} & \leanb{Development/\allowbreak{}Automorphic\allowbreak{}Arithmetic/\allowbreak{}Dedekind\allowbreak{}LFunctions/\allowbreak{}Char\allowbreak{}Family.\allowbreak{}lean:\allowbreak{}10-\allowbreak{}36}\\
\addlinespace[2pt]
Landau on windows & \textsf{S} & \leanb{\ldots{}/\allowbreak{}Dedekind\allowbreak{}LFunctions/\allowbreak{}Landau\allowbreak{}Theorems/\allowbreak{}Zero\allowbreak{}Free\allowbreak{}Of\allowbreak{}Rpow\allowbreak{}Window.\allowbreak{}lean:\allowbreak{}11}\\
\addlinespace[2pt]
Its proof: the window hypothesis in the logarithmic variable, zeros with real part at least 1 excluded separately, iterated derivatives of a Laplace transform under the bound C$\cdot$e\textasciicircum{}(2u), a kernel integral, a dichotomy and a nearest-zero lemma & \textsf{P} & \leanb{\ldots{}/\allowbreak{}Landau\allowbreak{}Theorems/\allowbreak{}lemma\allowbreak{}L.\allowbreak{}lean:\allowbreak{}13-\allowbreak{}23}; statements at \leanb{Iterated\allowbreak{}Deriv\allowbreak{}Laplace.\allowbreak{}lean:\allowbreak{}84}, \leanb{Landau\allowbreak{}Dichotomy.\allowbreak{}lean:\allowbreak{}24}, \leanb{Landau\allowbreak{}Nearest.\allowbreak{}lean:\allowbreak{}36}, \leanb{kernel\allowbreak{}Integral.\allowbreak{}lean:\allowbreak{}9}\\
\addlinespace[2pt]
A level and its function & \textsf{S} & \leanb{Level}, \leanb{Level.\allowbreak{}Gm}, \leanb{Formalism/\allowbreak{}Sharp\allowbreak{}Restriction/\allowbreak{}Ramification\allowbreak{}Frobenius/\allowbreak{}Totient\allowbreak{}Galois\allowbreak{}Classes/\allowbreak{}Explicit\allowbreak{}Formula\allowbreak{}Error\allowbreak{}Of\allowbreak{}Host\allowbreak{}Conditions.\allowbreak{}lean:\allowbreak{}8-\allowbreak{}39}; \leanb{KSet}, \leanb{\ldots{}/\allowbreak{}Totient\allowbreak{}Galois\allowbreak{}Classes/\allowbreak{}Prime\allowbreak{}Pair\allowbreak{}Size\allowbreak{}Condition.\allowbreak{}lean:\allowbreak{}39}\\
\addlinespace[2pt]
ASK, MAIN and the conditions on a host & \textsf{S} & \leanb{ASK}, \leanb{MAIN}, \leanb{De}, \leanb{Formalism/\allowbreak{}Size\allowbreak{}Window\allowbreak{}Brackets/\allowbreak{}First\allowbreak{}Prime\allowbreak{}Ne\allowbreak{}Second\allowbreak{}Prime.\allowbreak{}lean:\allowbreak{}84-\allowbreak{}93}; \leanb{Host\allowbreak{}Facts\allowbreak{}S}, \leanb{Development/\allowbreak{}Automorphic\allowbreak{}Arithmetic/\allowbreak{}Dedekind\allowbreak{}LFunctions/\allowbreak{}Split\allowbreak{}Inertia.\allowbreak{}lean:\allowbreak{}14}\\
\addlinespace[2pt]
The exact identity & \textsf{P} & \leanb{ch\_\allowbreak{}balance}, \leanb{\ldots{}/\allowbreak{}Dedekind\allowbreak{}LFunctions/\allowbreak{}Landau\allowbreak{}Theorems/\allowbreak{}Ch\allowbreak{}Balance.\allowbreak{}lean:\allowbreak{}200-\allowbreak{}262}\\
\addlinespace[2pt]
Thick and thin weights, kept characters, and the sums over them & \textsf{S} & \leanb{Sizes\allowbreak{}A}, \leanb{Chars\allowbreak{}A}, \leanb{Gm\allowbreak{}A}, \leanb{Sig\allowbreak{}A}, \leanb{gap\allowbreak{}A}, \leanb{\ldots{}/\allowbreak{}Dedekind\allowbreak{}LFunctions/\allowbreak{}Split\allowbreak{}Inertia.\allowbreak{}lean:\allowbreak{}26-\allowbreak{}66}; \leanb{Thin\allowbreak{}Scales}, \leanb{Formalism/\allowbreak{}Size\allowbreak{}Window\allowbreak{}Brackets/\allowbreak{}Growth\allowbreak{}Bound\allowbreak{}Gives\allowbreak{}Riemann\allowbreak{}Hypothesis.\allowbreak{}lean:\allowbreak{}55}; \leanb{Cost\allowbreak{}Windows}, \leanb{Formalism/\allowbreak{}Size\allowbreak{}Window\allowbreak{}Brackets/\allowbreak{}SChain.\allowbreak{}lean:\allowbreak{}50}\\
\addlinespace[2pt]
The scales of the thin places in the application & \textsf{S} & \leanb{Char\_\allowbreak{}at}, \leanb{Formalism/\allowbreak{}Sharp\allowbreak{}Restriction/\allowbreak{}Global\allowbreak{}Schur\allowbreak{}Bounds/\allowbreak{}Quadratic\allowbreak{}Log\allowbreak{}Deriv\allowbreak{}Congr.\allowbreak{}lean:\allowbreak{}221}\\
\addlinespace[2pt]
From the count to G, for kept characters & \textsf{P} & \leanb{ch\_\allowbreak{}core}, \leanb{\ldots{}/\allowbreak{}Landau\allowbreak{}Theorems/\allowbreak{}Ch\allowbreak{}Core.\allowbreak{}lean:\allowbreak{}343-\allowbreak{}444}; statements of \leanb{one\_\allowbreak{}char} and \leanb{core}, same file, 202 and 283\\
\addlinespace[2pt]
The bound for a level; the lattice inequality for each translate & \textsf{P} & \leanb{core\_\allowbreak{}gm}, \leanb{Development/\allowbreak{}Automorphic\allowbreak{}Arithmetic/\allowbreak{}Primitive\allowbreak{}Cosets/\allowbreak{}Core\allowbreak{}Gm.\allowbreak{}lean:\allowbreak{}12-\allowbreak{}74}; \leanb{ch\_\allowbreak{}top}, \leanb{\ldots{}/\allowbreak{}Landau\allowbreak{}Theorems/\allowbreak{}Ch\allowbreak{}Top.\allowbreak{}lean:\allowbreak{}37-\allowbreak{}74}; statements of \leanb{gm\allowbreak{}A\_\allowbreak{}eq}, \leanb{sig\allowbreak{}A\_\allowbreak{}le}, \leanb{EAt\allowbreak{}H}\\
\addlinespace[2pt]
What a host adapted to a level is; which kept characters are trivial on H & \textsf{S} & \leanb{EAt\allowbreak{}H}, \leanb{Formalism/\allowbreak{}Sharp\allowbreak{}Restriction/\allowbreak{}Ramification\allowbreak{}Frobenius/\allowbreak{}Totient\allowbreak{}Galois\allowbreak{}Classes/\allowbreak{}Explicit\allowbreak{}Formula\allowbreak{}Error\allowbreak{}Of\allowbreak{}Host\allowbreak{}Conditions.\allowbreak{}lean:\allowbreak{}144}; \leanb{Exceptional\allowbreak{}Character\allowbreak{}Kept}, \leanb{\ldots{}/\allowbreak{}Totient\allowbreak{}Galois\allowbreak{}Classes/\allowbreak{}Analytic.\allowbreak{}lean:\allowbreak{}158}; \leanb{wt\allowbreak{}A}, \leanb{\ldots{}/\allowbreak{}Dedekind\allowbreak{}LFunctions/\allowbreak{}Split\allowbreak{}Inertia.\allowbreak{}lean:\allowbreak{}32}; statement of \leanb{gm\allowbreak{}A\_\allowbreak{}eq\_\allowbreak{}triv}, \leanb{Development/\allowbreak{}Automorphic\allowbreak{}Arithmetic/\allowbreak{}Primitive\allowbreak{}Cosets/\allowbreak{}Gm\allowbreak{}AEq.\allowbreak{}lean:\allowbreak{}16}; \leanb{to\allowbreak{}E}, \leanb{\ldots{}/\allowbreak{}Primitive\allowbreak{}Cosets/\allowbreak{}Rec\allowbreak{}H.\allowbreak{}lean:\allowbreak{}10}\\
\addlinespace[2pt]
The size conditions and the three costs & \textsf{S} & \leanb{Size\allowbreak{}Windows}, \leanb{Growth\allowbreak{}Size\allowbreak{}Bound}, \leanb{Formalism/\allowbreak{}Size\allowbreak{}Window\allowbreak{}Brackets/\allowbreak{}Growth\allowbreak{}Bound\allowbreak{}Gives\allowbreak{}Riemann\allowbreak{}Hypothesis.\allowbreak{}lean:\allowbreak{}109-\allowbreak{}123}; \leanb{Rest\allowbreak{}Size}, \leanb{SChain.\allowbreak{}lean:\allowbreak{}53}\\
\addlinespace[2pt]
The record of eleven clauses & \textsf{S} & \leanb{Size\allowbreak{}Data\textsubscript{2}}, \leanb{Formalism/\allowbreak{}Size\allowbreak{}Window\allowbreak{}Brackets/\allowbreak{}SChain.\allowbreak{}lean:\allowbreak{}17}\\
\addlinespace[2pt]
The space of forms is all algebraic modular forms of the weight and level & \textsf{S} & \leanb{Fi}, \leanb{Development/\allowbreak{}Abelian\allowbreak{}Variety\allowbreak{}Theory/\allowbreak{}Definitions/\allowbreak{}Automorphic\allowbreak{}Membership/\allowbreak{}Pred\allowbreak{}Mem\allowbreak{}Gam\allowbreak{}Ranks.\allowbreak{}lean:\allowbreak{}151}\\
\addlinespace[2pt]
The rescaled lattice & \textsf{S} & \leanb{Iu}, \leanb{Nprime\allowbreak{}Def}, \leanb{Formalism/\allowbreak{}Local\allowbreak{}Ring\allowbreak{}Cardinality/\allowbreak{}Data\allowbreak{}DIVDefs.\allowbreak{}lean:\allowbreak{}17-\allowbreak{}24}\\
\addlinespace[2pt]
The clauses Loss, Gain, the comparison of Ee and the count for T & \textsf{S} & \leanb{Loss\allowbreak{}Bounds}, \leanb{Quotient\allowbreak{}Slope\allowbreak{}Bounds\allowbreak{}With}, \leanb{EComp5}, \leanb{Count5}, \leanb{Formalism/\allowbreak{}Size\allowbreak{}Window\allowbreak{}Brackets/\allowbreak{}Quotient\allowbreak{}Slope\allowbreak{}Bounds\allowbreak{}With.\allowbreak{}lean:\allowbreak{}19-\allowbreak{}35}\\
\addlinespace[2pt]
The clauses Slope and Div; the clauses that bound the slopes of the lattices of forms and the variation of the heights & \textsf{S} & \leanb{Slope\allowbreak{}Ineq}, \leanb{DIV}, \leanb{Slopes}, \leanb{HComp}, \leanb{Formalism/\allowbreak{}Rank\allowbreak{}Gap\allowbreak{}Bounds/\allowbreak{}Mean\allowbreak{}Bound\allowbreak{}Second\allowbreak{}Form.\allowbreak{}lean:\allowbreak{}118-\allowbreak{}156}\\
\addlinespace[2pt]
Periods of the unscaled lattice are divisible to the order of the rescaling exponent & \textsf{S} & \leanb{adm2}, \leanb{Formalism/\allowbreak{}Sharp\allowbreak{}Restriction/\allowbreak{}Divisibility/\allowbreak{}Local\allowbreak{}Quadratic\allowbreak{}Parts.\allowbreak{}lean:\allowbreak{}206}\\
\addlinespace[2pt]
The gain formula: positions along the Frobenius orbit, partners, capacities, excluded exponents, candidate primes & \textsf{S} & \leanb{Is\allowbreak{}V}, \leanb{Good}, \leanb{cap\allowbreak{}U}, \leanb{Excl}, \leanb{g0}, \leanb{g\allowbreak{}U}, \leanb{Cand}, \leanb{VV}, \leanb{gain\allowbreak{}P}, \leanb{Formalism/\allowbreak{}Sharp\allowbreak{}Restriction/\allowbreak{}Growth\allowbreak{}Bounds/\allowbreak{}Is\allowbreak{}Residue\allowbreak{}Class\allowbreak{}Zero.\allowbreak{}lean:\allowbreak{}179-\allowbreak{}361}; \leanb{sg}, \leanb{Fo}, \leanb{Ga}, \leanb{Ee}, \leanb{Formalism/\allowbreak{}Local\allowbreak{}Ring\allowbreak{}Cardinality/\allowbreak{}Data\allowbreak{}Params\textsubscript{2}.\allowbreak{}lean:\allowbreak{}212-\allowbreak{}224}\\
\addlinespace[2pt]
The slope inequality for nearly tight frames & \textsf{S} & \leanb{Theorem\allowbreak{}F}, \leanb{Formalism/\allowbreak{}Harder\allowbreak{}Narasimhan/\allowbreak{}Statement\allowbreak{}Frame\allowbreak{}F.\allowbreak{}lean}\\
\addlinespace[2pt]
Its proof & \textsf{P} & \leanb{theorem\allowbreak{}FSharp\_\allowbreak{}of}, \leanb{Formalism/\allowbreak{}Harder\allowbreak{}Narasimhan/\allowbreak{}Harder\allowbreak{}Narasimhan\allowbreak{}Proofs/\allowbreak{}Theorem\allowbreak{}FSharp.\allowbreak{}lean:\allowbreak{}168-\allowbreak{}390}; \leanb{frame\allowbreak{}B\_\allowbreak{}share}, \leanb{Frame\allowbreak{}Herm\allowbreak{}Sharp\allowbreak{}B.\allowbreak{}lean:\allowbreak{}25}\\
\addlinespace[2pt]
The ledger over the lift steps & \textsf{P} & \leanb{divisor\_\allowbreak{}inequality\_\allowbreak{}assembly}, \leanb{Formalism/\allowbreak{}Quotient\allowbreak{}Metric\allowbreak{}Divisors/\allowbreak{}Hermitian\allowbreak{}Twist\allowbreak{}Divisors/\allowbreak{}Div\allowbreak{}Assembly.\allowbreak{}lean:\allowbreak{}11-\allowbreak{}46}\\
\addlinespace[2pt]
The inequality for slopes under a factorisation of determinants & \textsf{S} & \leanb{Bounds\allowbreak{}Mu}, \leanb{Formalism/\allowbreak{}Divisor\allowbreak{}Inequality/\allowbreak{}Statement\allowbreak{}Obj.\allowbreak{}lean:\allowbreak{}17}\\
\addlinespace[2pt]
The lower bound for the scalars of the lifts and its seven hypotheses & \textsf{S} & \leanb{lam\allowbreak{}Low\_\allowbreak{}of\_\allowbreak{}hypotheses\_\allowbreak{}alt\_\allowbreak{}alt}, \leanb{\ldots{}/\allowbreak{}Third\allowbreak{}Place\allowbreak{}Series/\allowbreak{}Lam\allowbreak{}Low\allowbreak{}Of\allowbreak{}Published.\allowbreak{}lean:\allowbreak{}17}\\
\addlinespace[2pt]
The auxiliary lattices: integral forms on a moduli space & \textsf{S} & \leanb{cusp\allowbreak{}A}, \leanb{nrm\allowbreak{}M\_\allowbreak{}le\_\allowbreak{}one\_\allowbreak{}iff}, \leanb{Formalism/\allowbreak{}Local\allowbreak{}To\allowbreak{}Global/\allowbreak{}Projective\allowbreak{}Lattices/\allowbreak{}Cuspidal\allowbreak{}L2\allowbreak{}Subspace.\allowbreak{}lean:\allowbreak{}47-\allowbreak{}60}\\
\addlinespace[2pt]
The covolume clause and the chain of slope comparisons & \textsf{S} & \leanb{L3}, \leanb{Formalism/\allowbreak{}Local\allowbreak{}To\allowbreak{}Global/\allowbreak{}Indefinite\allowbreak{}Forms/\allowbreak{}Indef\allowbreak{}Side.\allowbreak{}lean:\allowbreak{}315}; \leanb{lcv\_\allowbreak{}le}, \leanb{Development/\allowbreak{}Abelian\allowbreak{}Variety\allowbreak{}Theory/\allowbreak{}Proofs/\allowbreak{}Divisor\allowbreak{}Proofs/\allowbreak{}Covolume\allowbreak{}From\allowbreak{}Gram\allowbreak{}And\allowbreak{}Bounds.\allowbreak{}lean:\allowbreak{}39}; its inputs at \leanb{Development/\allowbreak{}Abelian\allowbreak{}Variety\allowbreak{}Theory/\allowbreak{}Definitions/\allowbreak{}Divisors/\allowbreak{}Covol\allowbreak{}XOf\allowbreak{}Covol.\allowbreak{}lean:\allowbreak{}53-\allowbreak{}103}\\
\addlinespace[2pt]
The hypotheses of the covolume theorem from arithmetic geometry & \textsf{S} & \leanb{top\allowbreak{}Covol\allowbreak{}Slim2\_\allowbreak{}use\_\allowbreak{}alt\_\allowbreak{}alt}, \leanb{\ldots{}/\allowbreak{}Third\allowbreak{}Place\allowbreak{}Series/\allowbreak{}Top\allowbreak{}Covol\allowbreak{}Slim2\allowbreak{}Use.\allowbreak{}lean:\allowbreak{}16}\\
\addlinespace[2pt]
The Chowla--Selberg formula and Kronecker's limit formula & \textsf{S} & \leanb{chowla\_\allowbreak{}selberg\_\allowbreak{}eta} with the hypotheses of its section, \leanb{Development/\allowbreak{}Abelian\allowbreak{}Variety\allowbreak{}Theory/\allowbreak{}Proofs/\allowbreak{}Homogeneous\allowbreak{}Ideals/\allowbreak{}Dedekind\allowbreak{}Zeta\allowbreak{}Eq\allowbreak{}Sum\allowbreak{}Real\allowbreak{}Eisenstein.\allowbreak{}lean:\allowbreak{}212-\allowbreak{}225}; \leanb{tendsto\_\allowbreak{}real\allowbreak{}Eisenstein\_\allowbreak{}sub\_\allowbreak{}pole}, \leanb{\ldots{}/\allowbreak{}Tendsto\allowbreak{}Real\allowbreak{}Eisenstein\allowbreak{}Sub\allowbreak{}Pole.\allowbreak{}lean:\allowbreak{}352}\\
\addlinespace[2pt]
The class number formula and the level prime & \textsf{S} & \leanb{Class\allowbreak{}Number\allowbreak{}Formula\textsubscript{2}}, \leanb{Formalism/\allowbreak{}Sharp\allowbreak{}Restriction/\allowbreak{}Global\allowbreak{}Schur\allowbreak{}Bounds/\allowbreak{}Standard\allowbreak{}Lattice.\allowbreak{}lean:\allowbreak{}453}; \leanb{Level\allowbreak{}Prime}, \leanb{Development/\allowbreak{}Abelian\allowbreak{}Variety\allowbreak{}Theory/\allowbreak{}Definitions/\allowbreak{}Global\allowbreak{}Twisting/\allowbreak{}Log\allowbreak{}Cond\allowbreak{}Le\allowbreak{}Three.\allowbreak{}lean:\allowbreak{}81}\\
\addlinespace[2pt]
The lattice side equals MAIN identically & \textsf{S} & \leanb{key\_\allowbreak{}identity5}, \leanb{Formalism/\allowbreak{}Size\allowbreak{}Window\allowbreak{}Brackets/\allowbreak{}Chain\allowbreak{}Slope\allowbreak{}Windows/\allowbreak{}Level\allowbreak{}Of\allowbreak{}Congruence.\allowbreak{}lean:\allowbreak{}105}\\
\addlinespace[2pt]
The height-two logarithm; the radius lemma; the general form and the rate & \textsf{S} & \leanb{lt\allowbreak{}Log}, \leanb{Formalism/\allowbreak{}Local\allowbreak{}To\allowbreak{}Global/\allowbreak{}Torsion\allowbreak{}Integrality/\allowbreak{}CMGood\allowbreak{}Models/\allowbreak{}Statement.\allowbreak{}lean:\allowbreak{}35}; \leanb{Lemma\allowbreak{}St}, \leanb{Formalism/\allowbreak{}Gauss\allowbreak{}Jet\allowbreak{}Bounds/\allowbreak{}Statement.\allowbreak{}lean:\allowbreak{}107}; \leanb{Lemma\allowbreak{}RA}, \leanb{Formalism/\allowbreak{}Gauss\allowbreak{}Jet\allowbreak{}Bounds/\allowbreak{}Statement\allowbreak{}G.\allowbreak{}lean:\allowbreak{}12}\\
\addlinespace[2pt]
Good models of CM curves and their formal groups at inert primes & \textsf{S} & \leanb{CM2}, \leanb{CM3}, \leanb{\ldots{}/\allowbreak{}CMGood\allowbreak{}Models/\allowbreak{}Statement.\allowbreak{}lean:\allowbreak{}45-\allowbreak{}55}\\
\addlinespace[2pt]
Good reduction is derived from Artin reciprocity & \textsf{S} & \leanb{cm\allowbreak{}Good\allowbreak{}Model\_\allowbreak{}of\_\allowbreak{}global\allowbreak{}Artin\allowbreak{}Reciprocity}, \leanb{Development/\allowbreak{}Abelian\allowbreak{}Variety\allowbreak{}Theory/\allowbreak{}Proofs/\allowbreak{}Decorated\allowbreak{}Lattice\allowbreak{}Proofs/\allowbreak{}Cm\allowbreak{}Good\allowbreak{}Model\allowbreak{}Of\allowbreak{}Global\allowbreak{}Artin\allowbreak{}Reciprocity.\allowbreak{}lean:\allowbreak{}10}\\
\addlinespace[2pt]
Integrality of the theta kernel at inert primes & \textsf{S} & \leanb{Int\allowbreak{}Inert}, \leanb{Formalism/\allowbreak{}Local\allowbreak{}To\allowbreak{}Global/\allowbreak{}Inert\allowbreak{}Kernel\allowbreak{}Theory/\allowbreak{}Statement\allowbreak{}Int.\allowbreak{}lean:\allowbreak{}18}\\
\addlinespace[2pt]
The jet theorem is applied to an arbitrary series with the integrality hypothesis & \textsf{P} & \leanb{prv\allowbreak{}F\_\allowbreak{}of\_\allowbreak{}orbit\allowbreak{}Letter\allowbreak{}Bounds}, \leanb{Formalism/\allowbreak{}Period\allowbreak{}Integrals/\allowbreak{}Unitary\allowbreak{}Group\allowbreak{}Mass\allowbreak{}Formula/\allowbreak{}Orbit\allowbreak{}Valuation\allowbreak{}Bounds/\allowbreak{}Prv\allowbreak{}FOf\allowbreak{}Orbit\allowbreak{}Letter\allowbreak{}Bounds.\allowbreak{}lean:\allowbreak{}17-\allowbreak{}61}\\
\addlinespace[2pt]
Duality of the integral Weyl modules; Minkowski's bound & \textsf{S} & \leanb{standard\_\allowbreak{}weyl\_\allowbreak{}lattices\_\allowbreak{}fock\_\allowbreak{}dual}, \leanb{Development/\allowbreak{}Abelian\allowbreak{}Variety\allowbreak{}Theory/\allowbreak{}Proofs/\allowbreak{}Foundation\allowbreak{}Proofs/\allowbreak{}Standard\allowbreak{}Weyl\allowbreak{}Lattices\allowbreak{}Fock\allowbreak{}Dual\allowbreak{}Statement.\allowbreak{}lean:\allowbreak{}14}; \leanb{Minkowski\allowbreak{}Bound\allowbreak{}Q}, \leanb{Formalism/\allowbreak{}Period\allowbreak{}Integrals/\allowbreak{}Minkowski\allowbreak{}Q.\allowbreak{}lean:\allowbreak{}11}\\
\addlinespace[2pt]
The frame hypothesis & \textsf{S} & \leanb{Input\allowbreak{}SG}, \leanb{Formalism/\allowbreak{}Size\allowbreak{}Window\allowbreak{}Brackets/\allowbreak{}Frame\allowbreak{}Trace\allowbreak{}Ratio\allowbreak{}Fin\allowbreak{}Equiv.\allowbreak{}lean:\allowbreak{}30}\\
\addlinespace[2pt]
The period map is a Bessel pairing up to isometries and a non-zero scalar & \textsf{S} & \leanb{Bessel\allowbreak{}Id}, \leanb{Formalism/\allowbreak{}Sharp\allowbreak{}Restriction/\allowbreak{}Global\allowbreak{}Theorem\allowbreak{}Inputs.\allowbreak{}lean:\allowbreak{}65}\\
\addlinespace[2pt]
The expansion of the pairing; the identity term; vanishing for different indices & \textsf{S} & \leanb{KExpansion}, \leanb{Lemma\allowbreak{}V}, \leanb{Formalism/\allowbreak{}Local\allowbreak{}To\allowbreak{}Global/\allowbreak{}Siegel\allowbreak{}Lattice\allowbreak{}Theory/\allowbreak{}Closed\allowbreak{}Local\allowbreak{}Orbits/\allowbreak{}Term\allowbreak{}Op\allowbreak{}Stmt.\allowbreak{}lean:\allowbreak{}23-\allowbreak{}36}; \leanb{oot\_\allowbreak{}eq}, \leanb{\ldots{}/\allowbreak{}Closed\allowbreak{}Local\allowbreak{}Orbits/\allowbreak{}Gauge\allowbreak{}Lattice\allowbreak{}Shifts/\allowbreak{}Id\allowbreak{}Blocks.\allowbreak{}lean:\allowbreak{}672}; \leanb{Schur\allowbreak{}NFC}, \leanb{\ldots{}/\allowbreak{}Siegel\allowbreak{}Lattice\allowbreak{}Theory/\allowbreak{}Cokernel\allowbreak{}GLCounts/\allowbreak{}Rank\allowbreak{}Point\allowbreak{}Counts/\allowbreak{}Kappa\allowbreak{}Off\allowbreak{}Diagonal.\allowbreak{}lean:\allowbreak{}47}\\
\addlinespace[2pt]
Uniqueness of Bessel models with a gap & \textsf{S} & \leanb{Uniq\allowbreak{}Gap}, \leanb{Companion\allowbreak{}Theory/\allowbreak{}Automorphic\allowbreak{}Theory/\allowbreak{}Local\allowbreak{}Bessel\allowbreak{}Uniqueness/\allowbreak{}Bessel\allowbreak{}Models/\allowbreak{}At\allowbreak{}Form\allowbreak{}Adm\allowbreak{}One\allowbreak{}Of\allowbreak{}Steps.\allowbreak{}lean:\allowbreak{}15}\\
\addlinespace[2pt]
Siegel--Weil as a mass formula; the main row & \textsf{S} & \leanb{Siegel\allowbreak{}Densities}, \leanb{Formalism/\allowbreak{}Unipotent\allowbreak{}Adelic\allowbreak{}Quotients/\allowbreak{}Siegel\allowbreak{}Weil\allowbreak{}Unitary\allowbreak{}Convergent.\allowbreak{}lean:\allowbreak{}18}; \leanb{row\allowbreak{}Main\allowbreak{}Eq\allowbreak{}P\_\allowbreak{}of\_\allowbreak{}siegel\allowbreak{}Weil}, \leanb{Development/\allowbreak{}Abelian\allowbreak{}Variety\allowbreak{}Theory/\allowbreak{}Proofs/\allowbreak{}Flag\allowbreak{}Measures/\allowbreak{}Row\allowbreak{}Main\allowbreak{}Eq\allowbreak{}POf\allowbreak{}Siegel\allowbreak{}Weil.\allowbreak{}lean:\allowbreak{}10}\\
\addlinespace[2pt]
The bound for two different indices & \textsf{S} & \leanb{pb\allowbreak{}Applied}, \leanb{Formalism/\allowbreak{}Local\allowbreak{}To\allowbreak{}Global/\allowbreak{}Siegel\allowbreak{}Lattice\allowbreak{}Theory/\allowbreak{}Flag\allowbreak{}Stabilizer\allowbreak{}Tensors/\allowbreak{}Theorem\allowbreak{}S.\allowbreak{}lean:\allowbreak{}19}\\
\addlinespace[2pt]
In the application the parameter k of the frame is the rank parameter of the data & \textsf{P} & \leanb{Formalism/\allowbreak{}Local\allowbreak{}To\allowbreak{}Global/\allowbreak{}Siegel\allowbreak{}Lattice\allowbreak{}Theory/\allowbreak{}Flag\allowbreak{}Stabilizer\allowbreak{}Tensors/\allowbreak{}Theorem\allowbreak{}SClosed\allowbreak{}P.\allowbreak{}lean:\allowbreak{}102}\\
\addlinespace[2pt]
The cohomological bound and its windows & \textsf{S} & \leanb{VVle}, \leanb{VWin\allowbreak{}At}, \leanb{VWin}, \leanb{Development/\allowbreak{}Automorphic\allowbreak{}Arithmetic/\allowbreak{}Character\allowbreak{}Boxes/\allowbreak{}Exists\allowbreak{}Page\allowbreak{}N.\allowbreak{}lean:\allowbreak{}12-\allowbreak{}49}\\
\addlinespace[2pt]
The clause through which it is used, and the weakening of the square root & \textsf{P} & \leanb{Link\allowbreak{}N\_\allowbreak{}at}, \leanb{DIVin\allowbreak{}N\_\allowbreak{}at}, \leanb{Formalism/\allowbreak{}Sharp\allowbreak{}Restriction/\allowbreak{}Global\allowbreak{}Theorem\allowbreak{}Inputs.\allowbreak{}lean:\allowbreak{}31-\allowbreak{}63}; \leanb{divin\allowbreak{}N\_\allowbreak{}of\_\allowbreak{}hyp\allowbreak{}V}, \leanb{vv\allowbreak{}Record}, \leanb{Development/\allowbreak{}Automorphic\allowbreak{}Arithmetic/\allowbreak{}Character\allowbreak{}Boxes/\allowbreak{}Divin\allowbreak{}NOf\allowbreak{}Hyp\allowbreak{}V.\allowbreak{}lean:\allowbreak{}17-\allowbreak{}131}\\
\addlinespace[2pt]
The two quantities in which the bound is stated & \textsf{S} & \leanb{BNBounds}, \leanb{Formalism/\allowbreak{}Size\allowbreak{}Window\allowbreak{}Brackets/\allowbreak{}SChain.\allowbreak{}lean:\allowbreak{}33}\\
\addlinespace[2pt]
It is used in one place besides its producer & \textsf{S} & a search of the source for \leanb{DIVin\allowbreak{}N\_\allowbreak{}at}\\
\addlinespace[2pt]
The congruence ideal, the congruence module, and their bound & \textsf{S} & \leanb{cc\allowbreak{}Ideal}, \leanb{eta\allowbreak{}Loc}, \leanb{vcc}, \leanb{Formalism/\allowbreak{}Sharp\allowbreak{}Restriction/\allowbreak{}Level\allowbreak{}Letter\allowbreak{}Cocycles/\allowbreak{}Data\allowbreak{}N.\allowbreak{}lean:\allowbreak{}25-\allowbreak{}225}; \leanb{chi\_\allowbreak{}def}, same file, 126; \leanb{ZN}, same file, 1292; \leanb{Cong\allowbreak{}Bound\_\allowbreak{}at}, \leanb{Is\allowbreak{}Rank\allowbreak{}RPart}, \leanb{Development/\allowbreak{}Abelian\allowbreak{}Variety\allowbreak{}Theory/\allowbreak{}Block\allowbreak{}Theory/\allowbreak{}Non\allowbreak{}Galois\allowbreak{}Cusps/\allowbreak{}Base\allowbreak{}Sizes.\allowbreak{}lean:\allowbreak{}57-\allowbreak{}93}\\
\addlinespace[2pt]
The assembly of the bound & \textsf{S} & \leanb{of\_\allowbreak{}four\_\allowbreak{}alt}, \leanb{Development/\allowbreak{}Automorphic\allowbreak{}Arithmetic/\allowbreak{}Automorphic\allowbreak{}Varieties/\allowbreak{}Modularity\allowbreak{}Argument/\allowbreak{}Threshold\allowbreak{}Arguments/\allowbreak{}Dense\allowbreak{}Threshold\allowbreak{}Cases/\allowbreak{}HVCore.\allowbreak{}lean:\allowbreak{}18-\allowbreak{}22}\\
\addlinespace[2pt]
The crude bound for the second cohomology group & \textsf{S} & \leanb{Lemma\allowbreak{}K}, \leanb{Development/\allowbreak{}Automorphic\allowbreak{}Arithmetic/\allowbreak{}Automorphic\allowbreak{}Varieties/\allowbreak{}Generator\allowbreak{}Count\allowbreak{}Tuples/\allowbreak{}Lemma\allowbreak{}K.\allowbreak{}lean:\allowbreak{}13}\\
\addlinespace[2pt]
The odd part: from Poitou--Tate duality, with Kolyvagin derivative classes & \textsf{S} & \leanb{ML.\allowbreak{}iv}, \leanb{Development/\allowbreak{}Automorphic\allowbreak{}Arithmetic/\allowbreak{}Automorphic\allowbreak{}Varieties/\allowbreak{}Primitive\allowbreak{}Frobenius\allowbreak{}Fixed/\allowbreak{}Iv.\allowbreak{}lean:\allowbreak{}4}; \leanb{Psi\allowbreak{}Bound}, \leanb{\ldots{}/\allowbreak{}Automorphic\allowbreak{}Varieties/\allowbreak{}Odd\allowbreak{}Residue\allowbreak{}Bounds/\allowbreak{}Zero\allowbreak{}Mem\allowbreak{}Unr\allowbreak{}Hom.\allowbreak{}lean:\allowbreak{}97}; \leanb{kol\allowbreak{}Step}, \leanb{\ldots{}/\allowbreak{}Kol\allowbreak{}Step.\allowbreak{}lean:\allowbreak{}135}\\
\addlinespace[2pt]
Ferrero--Washington by Sinnott's theorem & \textsf{S} & \leanb{sinnott}, \leanb{Development/\allowbreak{}Automorphic\allowbreak{}Arithmetic/\allowbreak{}Automorphic\allowbreak{}Varieties/\allowbreak{}Galois\allowbreak{}Library/\allowbreak{}Ferrero\allowbreak{}Washington\allowbreak{}Index/\allowbreak{}Sinnott.\allowbreak{}lean:\allowbreak{}10}\\
\addlinespace[2pt]
The congruence number, the Eisenstein trace, and the bound of a length by valuations & \textsf{S} & \leanb{Number\allowbreak{}U3\allowbreak{}U\_\allowbreak{}s}, \leanb{One\allowbreak{}Char\allowbreak{}U\_\allowbreak{}s}, \leanb{Development/\allowbreak{}Automorphic\allowbreak{}Arithmetic/\allowbreak{}Automorphic\allowbreak{}Varieties/\allowbreak{}Modularity\allowbreak{}Argument/\allowbreak{}Threshold\allowbreak{}Arguments/\allowbreak{}Dense\allowbreak{}Threshold\allowbreak{}Cases/\allowbreak{}Number\allowbreak{}U3\allowbreak{}US.\allowbreak{}lean:\allowbreak{}30-\allowbreak{}60}; \leanb{eis3}, \leanb{\ldots{}/\allowbreak{}Automorphic\allowbreak{}Forms/\allowbreak{}Congruence\allowbreak{}Number\allowbreak{}Rank\allowbreak{}Three.\allowbreak{}lean:\allowbreak{}121}\\
\addlinespace[2pt]
The lemma on a graded local algebra & \textsf{S} & \leanb{Lemma\allowbreak{}M}, \leanb{Development/\allowbreak{}Automorphic\allowbreak{}Arithmetic/\allowbreak{}Automorphic\allowbreak{}Varieties/\allowbreak{}Commutative\allowbreak{}Algebra/\allowbreak{}Orbit\allowbreak{}Vanishing\allowbreak{}Local.\allowbreak{}lean:\allowbreak{}111}\\
\addlinespace[2pt]
The theorem on Hida families; Kato's Euler system & \textsf{S} & \leanb{hida\_\allowbreak{}unitary\_\allowbreak{}of\_\allowbreak{}printed\_\allowbreak{}hida}, \leanb{\ldots{}/\allowbreak{}Dense\allowbreak{}Threshold\allowbreak{}Cases/\allowbreak{}THida\allowbreak{}U.\allowbreak{}lean:\allowbreak{}38}; \leanb{zeta\allowbreak{}Ur}, \leanb{\ldots{}/\allowbreak{}Kato\allowbreak{}Quadratic\allowbreak{}Characters/\allowbreak{}Unramified\allowbreak{}Zeta/\allowbreak{}Zeta\allowbreak{}Ur.\allowbreak{}lean:\allowbreak{}9}\\
\addlinespace[2pt]
The partial compatibility of the Galois representations & \textsf{S} & \leanb{Away\allowbreak{}B}, \leanb{Development/\allowbreak{}Automorphic\allowbreak{}Arithmetic/\allowbreak{}Foundations/\allowbreak{}Galois\allowbreak{}Representations/\allowbreak{}Type\allowbreak{}Conditions/\allowbreak{}Weights\allowbreak{}Conj\allowbreak{}Bis.\allowbreak{}lean:\allowbreak{}26}\\
\addlinespace[2pt]
The theorem and its single hypothesis & \textsf{S} & \leanb{galois\allowbreak{}Representations\_\allowbreak{}of\_\allowbreak{}theta\allowbreak{}Lift}, \leanb{Development/\allowbreak{}Automorphic\allowbreak{}Arithmetic/\allowbreak{}Foundations/\allowbreak{}Galois\allowbreak{}Representations/\allowbreak{}Integral\allowbreak{}Structures/\allowbreak{}Galois\allowbreak{}Representations\allowbreak{}Of\allowbreak{}Theta\allowbreak{}Lift.\allowbreak{}lean:\allowbreak{}6}; \leanb{Row0b}, \leanb{\ldots{}/\allowbreak{}Type\allowbreak{}Conditions/\allowbreak{}Weights\allowbreak{}Conj\allowbreak{}Bis.\allowbreak{}lean:\allowbreak{}69}; \leanb{TLift\allowbreak{}K}, \leanb{\ldots{}/\allowbreak{}Galois\allowbreak{}Representations/\allowbreak{}Theta\allowbreak{}Parseval/\allowbreak{}Half\allowbreak{}Index\allowbreak{}Lattices/\allowbreak{}TLift\allowbreak{}K.\allowbreak{}lean:\allowbreak{}12}\\
\addlinespace[2pt]
The host field and the sizes of its primes & \textsf{S} & \leanb{Is\allowbreak{}Chi\allowbreak{}Field}, \leanb{Win\allowbreak{}Pair}, \leanb{Page\allowbreak{}N}, \leanb{Development/\allowbreak{}Automorphic\allowbreak{}Arithmetic/\allowbreak{}Character\allowbreak{}Boxes/\allowbreak{}Exists\allowbreak{}Page\allowbreak{}N.\allowbreak{}lean:\allowbreak{}19-\allowbreak{}40}\\
\addlinespace[2pt]
The host contains the imaginary quadratic field and the CM type fixes it & \textsf{S} & \leanb{Host\allowbreak{}Raw}, \leanb{Formalism/\allowbreak{}Sharp\allowbreak{}Restriction/\allowbreak{}Ramification\allowbreak{}Frobenius/\allowbreak{}Totient\allowbreak{}Galois\allowbreak{}Classes/\allowbreak{}Conj\allowbreak{}EMem\allowbreak{}Host\allowbreak{}FIff.\allowbreak{}lean:\allowbreak{}12}\\
\addlinespace[2pt]
Supply of hosts on windows & \textsf{S} & \leanb{supply\allowbreak{}Chi}, \leanb{\ldots{}/\allowbreak{}Character\allowbreak{}Boxes/\allowbreak{}Supply\allowbreak{}Chi.\allowbreak{}lean:\allowbreak{}8}\\
\addlinespace[2pt]
Good pairs of primes exist on a range of scales & \textsf{S} & \leanb{Pair\allowbreak{}Good}, \leanb{Pair\allowbreak{}Ex}, \leanb{\ldots{}/\allowbreak{}Character\allowbreak{}Boxes/\allowbreak{}Pair\allowbreak{}Good.\allowbreak{}lean:\allowbreak{}19-\allowbreak{}32}\\
\addlinespace[2pt]
The mean square and the zero-density estimate & \textsf{S} & \leanb{mean\allowbreak{}Sq\allowbreak{}K}, \leanb{Development/\allowbreak{}Automorphic\allowbreak{}Arithmetic/\allowbreak{}Mean\allowbreak{}Square\allowbreak{}Bounds/\allowbreak{}mean\allowbreak{}Sq\allowbreak{}K.\allowbreak{}lean:\allowbreak{}51}; \leanb{zero\_\allowbreak{}density\_\allowbreak{}box}, \leanb{Formalism/\allowbreak{}GL3\allowbreak{}Mean\allowbreak{}Squares/\allowbreak{}Genus\allowbreak{}Mean\allowbreak{}Asymptotics/\allowbreak{}Jacobi\allowbreak{}Character\allowbreak{}Sums/\allowbreak{}Count\allowbreak{}ZAll.\allowbreak{}lean:\allowbreak{}13}\\
\addlinespace[2pt]
The thick set without small relations; the statement Green & \textsf{S} & \leanb{No\allowbreak{}Rel}, \leanb{Formalism/\allowbreak{}Frobenius\allowbreak{}Orbits/\allowbreak{}Host\allowbreak{}Thick\allowbreak{}Exists.\allowbreak{}lean:\allowbreak{}13}; \leanb{no\allowbreak{}Rel\_\allowbreak{}image}, \leanb{host\allowbreak{}Thick\allowbreak{}Exists}, \leanb{Formalism/\allowbreak{}Frobenius\allowbreak{}Orbits/\allowbreak{}Frobenius\allowbreak{}Orbits\allowbreak{}Proofs/\allowbreak{}Thick\allowbreak{}Exists.\allowbreak{}lean:\allowbreak{}85-\allowbreak{}181}; \leanb{Is\allowbreak{}Prime\allowbreak{}Lab}, \leanb{Green}, \leanb{Development/\allowbreak{}Abelian\allowbreak{}Variety\allowbreak{}Theory/\allowbreak{}Definitions/\allowbreak{}Non\allowbreak{}Galois\allowbreak{}Descent/\allowbreak{}Label\allowbreak{}Datum.\allowbreak{}lean:\allowbreak{}91-\allowbreak{}112}\\
\addlinespace[2pt]
What each of the 186 literature statements asserts & \textsf{S} & the Lean statement of each row and, where it is a conjunction of named parts, the head of each part, at the path recorded for the row in \leanb{refs/\allowbreak{}rows-\allowbreak{}A.\allowbreak{}jsonl}, \leanb{rows-\allowbreak{}B.\allowbreak{}jsonl}, \leanb{rows-\allowbreak{}C.\allowbreak{}jsonl}; our descriptions are in \leanb{refs/\allowbreak{}rows-\allowbreak{}author.\allowbreak{}json}\\
\addlinespace[2pt]
Whether each is the published theorem in the generality of a source, a restricted form, or a different statement & \textsf{C} & the Lean statement of each row; the verdict and its reason are recorded in \leanb{refs/\allowbreak{}rows-\allowbreak{}scope-\allowbreak{}author.\allowbreak{}json} and printed in the table of the next appendix\\
\addlinespace[2pt]
The sources of each row & \textsf{R} & the reference research recorded in the same three files\\
\addlinespace[2pt]
What the final theorem does not use & \textsf{M} & \leanb{results/\allowbreak{}cut-\allowbreak{}result.\allowbreak{}txt}, \leanb{cut-\allowbreak{}in.\allowbreak{}txt.\allowbreak{}xz}, \leanb{cut-\allowbreak{}out.\allowbreak{}txt.\allowbreak{}xz}, \leanb{cut-\allowbreak{}modules.\allowbreak{}txt.\allowbreak{}xz}; statements of the named declarations read at \leanb{Formalism/\allowbreak{}Product\allowbreak{}Formula/\allowbreak{}Statement\allowbreak{}OCCM.\allowbreak{}lean:\allowbreak{}137-\allowbreak{}160}, \leanb{Formalism/\allowbreak{}Half\allowbreak{}Links/\allowbreak{}Statement.\allowbreak{}lean:\allowbreak{}144}, \leanb{Formalism/\allowbreak{}Jacobi\allowbreak{}Weight\allowbreak{}Families/\allowbreak{}Statement.\allowbreak{}lean:\allowbreak{}57}, \leanb{Formalism/\allowbreak{}Statement\allowbreak{}OC.\allowbreak{}lean:\allowbreak{}65}, \leanb{Formalism/\allowbreak{}Size\allowbreak{}Window\allowbreak{}Brackets/\allowbreak{}First\allowbreak{}Prime\allowbreak{}Ne\allowbreak{}Second\allowbreak{}Prime.\allowbreak{}lean:\allowbreak{}47}\\
\addlinespace[2pt]
\end{longtable}}
